\input plain
\hoffset=0.5in
\voffset=0.5in
\hsize=5.5in
\vsize=9.5in
\parskip=6pt
\parindent=18pt

% Fonts
\font\titlefont=cmr12 scaled\magstep2
\font\sectionfont=cmbx12
\font\bodyfont=cmr12
\font\itfont=cmti12
\font\smallfont=cmr10

\bodyfont

% Centred display macro
\def\centrepar#1{\medskip\centerline{#1}\medskip}

% Section heading macro
\def\section#1{\medskip{\sectionfont #1}\smallskip}

%%%

\centerline{\titlefont An Information-Theoretic Reading of the List Type}
\smallskip
\centerline{\itfont A Pearl}
\bigskip

\section{The List Type as a Geometric Series}

A list over a type~$A$ may be defined by the union of all finite products:
$$\hbox{List}(A) \;=\; 1 + A + A^2 + A^3 + \cdots$$
This is a geometric series in~$A$, and by the standard identity it contracts to the closed form:
$$\hbox{List}(A) \;=\; {1 \over 1 - A}$$
This is not merely a formal trick. Each step of the equivalence carries meaning,
and reading it carefully~--- attending to how each symbol is naturally pronounced~---
yields a small constellation of insights.

\section{Pronouncing the Operators}

The operators of type algebra have natural readings. Addition~($+$) is {\itfont or\/}:
a disjoint choice between alternatives. Multiplication~($\times$) is {\itfont and\/}:
the simultaneous holding of two things. Their respective identities follow:
$0$, the empty type, is {\itfont impossibly\/}~--- the choice that offers nothing;
and $1$, the unit type, is {\itfont trivially\/}~--- the conjunction that adds no information.

Inverses must invert. Division~($/$) is therefore {\itfont regardless of\/}~---
the dissolution of a simultaneous distinction (as in $\mathbf{R}/\mathbf{Z}$:
the reals, regardless of integer part). And subtraction~($-$) inverts disjunction:
{\itfont without distinguishing\/}, or {\itfont pretending indifference between\/}.

Under these readings, $1/(1-A)$ becomes: {\itfont the trivial, regardless of
pretending-$A$-indistinguishable-from-the-trivial.\/} Equated to the series~---
``a list is nothing, or one thing, or two things,~$\ldots$''~--- the tautology
crystallises: the whole is nothing but the iterated overcoming of its own lack.

\section{The Radius of Convergence as a Logical Boundary}

The geometric series converges only when $|A| < 1$. This analytic condition is not
imported from complex analysis by accident~--- it is the condition that $A$ carry
genuine information. To construct a random instance of a type is to perform a
probabilistic event; $|A|$ is therefore literally the probability of that event,
constrained by construction to $[0,1]$.

The boundary cases are instructive. When $|A| = 0$, $A$ is the empty type~---
{\itfont impossibly\/}~--- and $\hbox{List}(A) = 1$: there is exactly one list of
impossible things, the empty list. Vacuous truth is well-behaved. When $|A| = 1$,
$A$ is the unit type~--- {\itfont trivially\/}~--- and the series diverges:
$1 + 1 + 1 + \cdots$. A list of tautologies is not a list; it is an explosion
of undifferentiated sameness.

Shannon makes this precise. The information content of an event of probability~$p$
is $-\log p$. The convergence condition $|A| < 1$ is therefore precisely
$-\log|A| > 0$: $A$ must carry positive information. The series~--- and with it
the well-foundedness of the list type~--- converges if and only if each element
says something.

In base~2 this becomes especially concrete. An element of probability $|A| = 1/2$
carries exactly one bit, corresponding to a fair Boolean coin-toss; and
$\hbox{List}(A)$ converges to~$2$, the two-element type. The familiar and the
foundational are the same object, viewed from different distances.

The analytic boundary and the logical boundary coincide: the list construction is
well-behaved precisely in the space between impossibility and triviality, between
$\bot$ and $\top$, between $0$ and~$1$. The radius of convergence was always a
logical condition dressed in analytic clothing. {\itfont The mathematics knew before we did.\/}

\section{The Exponential as Finite Sets}

The same algebra admits a richer construction. Observe that $X^n$ is an ordered
$n$-tuple: $n$ things held {\itfont with\/} their positional distinctions. The
symmetric group $S_n$, of order $n!$, acts on $X^n$ by permuting those positions;
quotienting by it~--- dividing, in our pronunciation: {\itfont regardless of
order\/}~--- yields the unordered $n$-element multiset $X^n/n!$. The factorial
is not a normalisation convenience; it {\itfont is\/} the symmetry group, and
division is the quotient by it.

The exponential power series therefore reads:
$$\exp(A) \;=\; \sum_{n=0}^\infty {A^n \over n!}
  \;=\; 1 + A + \{A,A\} + \{A,A,A\} + \cdots$$
which pronounces as: {\itfont trivially, or one thing, or an unordered pair, or
an unordered triple,~$\ldots$\/}~--- that is, a finite set over~$A$. The
exponential of a type is its type of finite sets. Every symbol is pronounceable;
nothing is opaque.

\section{Doubling, Periodicity, and Transcendence}

Passing from $\mathbf{R}$ to $\mathbf{C}$ via Cayley--Dickson doubling, we
evaluate the exponential at a purely imaginary argument $(0, y)$. The resulting
function is periodic, and its period is $2\pi$. Both the existence and the value
of this period are consequences of perfectly well-typed, pronounceable algebra.
Yet $2\pi$ is transcendental: it cannot be expressed as the root of any polynomial
with integer~--- that is, {\itfont counting\/}~--- coefficients.

This is the natural home of the notion of transcendence. The type algebra built
from $\{0, 1, {+}, {\times}\}$ and their inverses constitutes a language; its
closed terms name exactly the algebraic numbers. The integers are the
{\itfont counting types\/}~--- the canonical inhabitants of the arithmetic ground.
Algebraic numbers are precisely those nameable as roots of polynomials over this
ground: the values {\itfont expressible\/} in the language of type algebra.

Transcendental numbers are then the values that {\itfont exist\/}~--- that arise
as periods, or limits, of perfectly well-formed constructions in the language~---
but that {\itfont cannot be named\/} within it. Transcendence is incompleteness:
a truth the system can gesture toward but not utter. The period $2\pi$ is right
there, the function converges, the circle closes; yet no polynomial equation with
integer coefficients will ever pin it down.

This gives a precise and pronounceable meaning to the otherwise mysterious
designation {\itfont algebraic data type\/}: such types are algebraic in exactly
the sense that algebraic numbers are~--- constructible from the arithmetic ground
by the operations the algebra sanctions. The gap between algebraic and
transcendental numbers is the same gap as the gap between what a type algebra can
construct and what its own well-typed functions can witness but not name.

\section{A Type-Theoretic Gloss on Schanuel's Conjecture}

Schanuel's conjecture asserts that if $z_1, \ldots, z_n \in \mathbf{C}$ are
linearly independent over~$\mathbf{Q}$, then the $2n$ numbers
$z_1, \ldots, z_n, e^{z_1}, \ldots, e^{z_n}$ have transcendence degree at least
$n$ over~$\mathbf{Q}$. It implies, as special cases, the transcendence of $e$
and~$\pi$, their algebraic independence, and much else besides.

In the type-theoretic reading developed here, the conjecture has a natural gloss.
The rationals~$\mathbf{Q}$ are the ground of the type algebra: the counting
language, the coefficients that polynomials are permitted to use. Linear
independence over~$\mathbf{Q}$ is the condition that the $z_i$ carry genuinely
distinct information~--- that none is a rational linear combination, i.e.\ a
{\itfont type-algebraic expression\/}, of the others.

Schanuel's conjecture then says: applying the exponential~--- the finite-set
constructor, the passage from ordered to unordered, from dynamic to
geometric~--- to $n$ algebraically independent inputs produces outputs that are
collectively {\itfont maximally\/} inexpressible in the ground language. Each
application of $\exp$ contributes a full new dimension of transcendence; no
collapse occurs.

Viewed this way, the conjecture is a maximality axiom for the expressiveness gap:
it asserts that the exponential is {\itfont as transcendental as it possibly
could be\/}, given its inputs. And note the asymmetry of the two canonical
constants. $e = \exp(1)$ is the exponential evaluated at the multiplicative
unit~--- the dynamic, temporal fixed point of differentiation, the type of finite
sets over a single point. $\pi$ enters through the period of $\exp(0,y)$~--- the
geometric, spatial extremal invariant, the measure of the circle that
Cayley--Dickson doubling makes inevitable. Their algebraic independence is the
statement that the dynamic ground and the geometric ground are incommensurable
in the counting language: two species of transcendence, orthogonal to one another
and to everything the algebra can say.

Not many type theorists notice that {\itfont algebraic\/} data types are algebraic
in exactly the number-theoretic sense. Schanuel's conjecture suggests they are
right to have taken the word seriously: the boundary of the algebraic is where
the most important things live.

\medskip
\centerline{---}
\medskip

{\itfont That the same object~--- a humble list~--- simultaneously encodes a
tautology about lack and overcoming, a probabilistic constraint on
informativeness, the logical gap between the impossible and the trivial, and a
type-theoretic shadow of one of the deepest open problems in mathematics, is not
a coincidence. It is a small instance of the general fact that structure, once
found, ramifies everywhere. All is metaphor; the metaphors are load-bearing.\/}

\vfill
\bye